\section{Literature Synthesis and Open Challenges}
\label{sec:literature-synthesis}
\label{sec:literature-synthesis-open-challenges}

The preceding chapters introduced the review axes one at a time: continuum
fields, model hierarchy, closure choices, observables, numerical methods, and
evidence standards. This chapter compares across those axes. Its purpose is not
to add another catalog of sources, but to state what the landscape collectively
shows. The central finding is that blood-flow models are interpretable only
when model fidelity, retained state, data, closure choices, observables,
numerical evidence, and validation status are aligned with the question being
asked.

\subsection{Fidelity, Cost, and Retained State}
\label{subsec:synthesis-fidelity-identifiability-observables}

Increasing dimensional fidelity can expose local velocity gradients, wall
interaction, and geometry-specific flow features, but it also increases the
number of boundary, material, geometric, and numerical assumptions that must be
identified. Three-dimensional CFD and FSI formulations retain the most direct
representation of local flow and wall kinematics, while 1D and 0D models
deliberately compress that information into axial or lumped variables
\cite{FormaggiaEtAl2007FSICoupling,QuarteroniVeneziani2003GeometricalMultiscale,ShiEtAl2011BloodFlowModelReview}.
The reduced tiers are not weaker by definition. They are appropriate when the
question concerns pulse propagation, network pressure--flow behavior, rapid
screening, or parameter studies for which fully resolved data and boundary
conditions are unavailable or unnecessary.

The cost of that compression is interpretive. A 1D model may have fewer state
variables than a 3D model, but it can still be underdetermined if wall, profile,
friction, rheology, and boundary closures are treated as adjustable without
independent evidence. Conversely, a highly resolved forward model can depend on
material properties, prestress, inlet histories, outflow impedances, geometry
state, and wall motion that are rarely all observed with matching certainty.
Parameter-identifiability work for hemodynamics PDE models makes this tension
explicit: adding modeling detail is valuable only when the data and outputs can
distinguish the resulting parameters
\cite{Colebank2025HemodynamicsIdentifiability}.

\subsection{Observables and Cross-Model Comparability}
\label{subsec:synthesis-observables-comparability}

Observation operators are part of the model, not post-processing decoration. A
1D area--flow solution does not natively contain the same observable as a
resolved 3D velocity field, and a pressure-drop, flow-rate, section-mean
velocity, wall-shear, or FFR-style quantity each imposes a different comparison
contract. Multiscale coupling work already treats interface quantities and
compatibility conditions as mathematical structure
\cite{FormaggiaEtAl2007FSICoupling,QuarteroniVeneziani2003GeometricalMultiscale};
diagnostic 1D--3D comparison needs the same discipline at the output side.

The clinical FFR literature sharpens the point because it separates anatomical
stenosis appearance from pressure-flow function
\cite{ToninoEtAl2010FAMEAngioFunctional,TebaldiEtAl2015FFROverview}. CT-derived
FFR and 1D--3D FFR comparison sources agree that the diagnostic scalar is not a
raw geometry measurement; it is an output of an acquisition, modeling, boundary,
and pressure-observation chain
\cite{NorgaardEtAl2014NXTFFRCT,Blanco2018FFR1D3D,ShaikhEtAl2025CTFFR}. Thus a
cross-model discrepancy is uninterpretable until the observed quantity is
declared and matched.

\subsection{Uncertainty in Boundaries, Walls, Rheology, and Geometry}
\label{subsec:synthesis-closure-geometry-uncertainty}

Boundary-condition uncertainty is one of the most persistent barriers to
interpretable hemodynamics. Inlet waveforms, outlet impedances, terminal
resistances, and coupled-domain histories can change the local solution as much
as the internal discretization. In a multiscale setting, these quantities are
not auxiliary inputs; they are the mechanism by which upstream and downstream
vascular beds communicate with the modeled region
\cite{QuarteroniVeneziani2003GeometricalMultiscale,FormaggiaEtAl2007FSICoupling}.
A reduced stenosis result that does not declare its boundary approximation
therefore cannot be cleanly compared with a resolved result, even if both
models share a nominal geometry.

Wall and rheology uncertainty create a second closure layer. Wall stiffness and
thickness determine wave speed and area response; rheology and velocity-profile
assumptions determine dissipative behavior and momentum transport. In 1D
models, these effects enter through pressure--area laws, friction terms,
effective viscosity choices, and momentum-flux correction factors
\cite{ShiEtAl2011BloodFlowModelReview}. In FSI models, they enter through wall
material laws, structural boundary conditions, and the coupling algorithm. Both
settings face the same interpretive problem: when the wall or rheology is not
independently constrained, a discrepancy can be a physics signal, a closure
signal, or a compensation between uncertain choices.

Geometry is similarly double-edged. Patient-specific geometry is attractive
because it carries clinically meaningful morphology, but segmentation,
registration, current-versus-reference configuration, and wall-motion state can
dominate the modeling record. Idealized geometry removes many of those
ambiguities and exposes controlled severity and refinement parameters, but it
cannot by itself answer patient-specific questions. The useful distinction is
not realistic versus artificial. It is whether the geometry serves the claim:
patient-specific geometry for patient-specific interpretation, and idealized
geometry for isolating numerical, closure, and comparison behavior.

\subsection{Validation Status, Reproducibility, and Evidence Limits}
\label{subsec:synthesis-consistency-verification-validation}

Verification, validation, and reproducibility should not be collapsed.
Manufactured solutions and grid studies ask whether an implementation behaves
as expected for a declared mathematical problem. Benchmarks and open solvers
support reproducibility and comparison across workflows
\cite{BoileauEtAl2015Benchmark,BenemeritoEtAl2024OpenBF}. Resolved-flow
platforms and model repositories extend that reproducibility lesson to 3D
workflows by making geometry, meshing, solver, and repository records easier to
inspect
\cite{SimVascular2026Platform,SimVascular2026VascularModelRepository}. None of
these records is automatically a validation claim unless the model, data,
observable, and accepted reference match the target question.

Equilibrium preservation is a particularly sharp consistency test. A continuous
stenotic balance law may admit a rest or steady state through cancellation
between flux gradients, geometry terms, wall response, friction, and boundary
data. A discrete method that does not preserve that state can generate motion
from the numerical representation itself. Well-balanced methods for 1D blood
flow make this point directly and raise the standard for stenosis computations:
manufactured solutions and grid studies are useful, but an equilibrium problem
needs an equilibrium-preservation audit
\cite{PimentelGarciaEtAl2025HighOrderWellBalanced}. The broader verification
and validation gap is that many studies report plausible hemodynamic outputs
without enough evidence to separate implementation error, discretization error,
closure uncertainty, cross-model mismatch, and external validation against an
accepted target.

\subsection{Reproducibility and Emerging Learned Methods}
\label{subsec:synthesis-reproducibility-learned-methods}

Reproducibility is the practical bridge between modeling ambition and reviewable
evidence. Benchmark studies help by fixing common cases and comparison
quantities, but stenosis studies also need persisted geometry definitions,
boundary histories, closure parameters, solver settings, diagnostics, and output
maps \cite{BoileauEtAl2015Benchmark,BenemeritoEtAl2024OpenBF}. Without those
records, a disagreement between two models is hard to assign to physics,
parameters, numerics, or data handling.

Learned methods do not remove this requirement. Physics-informed neural networks
offer a framework for solving and inverting PDE-constrained problems
\cite{RaissiEtAl2019PhysicsInformedNeuralNetworks}, and hemodynamics reviews
now treat PINNs as an active direction for vascular modeling
\cite{YuEtAl2026PINNHemodynamicsReview}. At the same time, gradient-pathology
analyses show that training dynamics can be a limiting numerical issue rather
than a transparent solver detail
\cite{WangEtAl2021PINNGradientPathologies}. Operator-learning methods such as
DeepONet and Fourier neural operators shift the ambition from one solution to a
map between function spaces \cite{LuEtAl2021DeepONet,LiEtAl2021FNO}, and recent
hemodynamics-specific neural differential and stenosed-vessel operator-learning
work shows why such models are appealing for fast parameterized flow prediction
\cite{CsalaEtAl2025PCNDE,VelikorodnyEtAl2025DeepOperatorStenosed}. Their open
challenge is the same synthesis problem in a new form: the training
distribution, geometry family, boundary inputs, wall and rheology assumptions,
loss normalization, and validation observables must be declared before speed or
fit can be interpreted as hemodynamic evidence.

\subsection{Role of the Idealized Stenosis Case Study}
\label{subsec:synthesis-idealized-stenosis-case-study}

The idealized stenosis case study is useful because it applies the review
framework in one controlled setting. It fixes a geometry family, closure
contract, boundary approximation, discretization, and observation operator
tightly enough that the consequences can be inspected. The same setup can ask
whether a reduced model preserves a rest equilibrium, whether a velocity
comparison is being made through a declared cross-sectional operator, whether
the recorded artifacts are sufficient to reproduce the numerical claim, and
which uncertainties remain outside the case.

The case study is therefore not the destination of the literature review. It is
one auditable example of the review standard. It removes segmentation and
clinical-matching ambiguity so that closure dependence, boundary sensitivity,
equilibrium preservation, verification evidence, and cross-dimensional
observability can be examined in a controlled setting. Patient-specific
interpretation remains the domain of studies with matched patient geometry,
boundary data, wall state, and validation targets.
